% Source: Wald GR, physical PDF p. 214
%         (printed p. 201).
% Coverage: A future domain of dependence whose Cauchy horizon intersects \(S\).
% Figures/tables/footnotes: Figure 8.12; no tables or footnotes.
% Status: source-reviewed and compile-clean.
% Uncertainties: none.

\begingroup
\renewcommand{\thefigure}{8.12}
\begin{figure}[H]
  \centering
  \begin{tikzpicture}[x=.83cm,y=.83cm,>=Latex,line cap=round,font=\small]
    \draw[thick] (-3.15,-1.28) -- (-1.06,2.25) -- (2.98,-1.02);
    \draw[thick] (-3.15,-1.28) .. controls (-1.66,1.02) and (-.62,.87) .. (1.17,.15)
      .. controls (1.89,-.17) and (2.40,-.67) .. (2.98,-1.02);
    \node at (-.65,1.16) {\(D^+(S)\)};
    \draw[->] (-1.36,2.58) .. controls (-1.68,2.16) and (-1.75,1.89) .. (-1.78,1.62);
    \draw[->] (.40,2.52) .. controls (.36,2.05) and (.52,1.78) .. (.73,1.48);
    \node at (-.66,2.61) {\(H^+(S)\)};
    \draw[->] (.01,.03) -- (.16,.30) node[below=.10cm] {\(S\)};
  \end{tikzpicture}
  \caption{A spacetime diagram showing \(D^+(S)\) and \(H^+(S)\) for a particular
  closed achronal set \(S\) in Minkowski spacetime. Here \(S\) is \lq\lq asymptotically
  null\rq\rq\ to the right and becomes \lq\lq exactly null\rq\rq\ to the left. This example
  shows that \(H^+(S)\) can intersect \(S\) even if \(\operatorname{edge}(S)=\emptyset\).}
  \label{fig:8.12}
\end{figure}
\endgroup
